% ============================================================================
%  Actividad 5 - Conclusiones generales, referencias consolidadas y Anexo A
%
%  Este archivo cierra el documento acumulado: no cierra solo la prueba de
%  usabilidad, sino el caso de estudio completo A1-A5. Tres piezas:
%
%   1. Conclusiones generales: el arco del proyecto, el cierre del circulo
%      -que predijeron las Actividades 1 a 4 y que encontro la prueba del
%      20 de agosto de 2026- y el orden de la siguiente iteracion (R1-R7).
%   2. Referencias consolidadas del documento acumulado, en APA y con
%      sangria francesa (entorno uxreferencias de estilo/uxdoc.sty).
%   3. Anexo A con los instrumentos aplicados. El manual de uso es el Anexo B
%      y vive en contenido/a5_07_manual.tex: aqui no se escribe ni se incluye.
%
%  Reglas que aplican aqui: ninguna cifra sin origen -toda proviene de las
%  tablas de la seccion de prueba de usabilidad o de una cita-, ningun tiempo
%  futuro y ninguna autoria inventada: los trabajos de 2026 cuyos autores no
%  pudieron verificarse se citan por titulo e identificador arXiv.
%
%  La direccion del demo entra por la macro \urlDemoWeb de datos/demo.tex.
%  Nunca una URL literal en este archivo.
% ============================================================================

\clearpage
\section{Conclusiones generales}\label{sec:a5-conclusiones}

El proyecto recorrió cinco entregas y terminó en un punto que no estaba garantizado al empezar:
una experiencia que puede recorrerse, medirse y corregirse con evidencia de personas reales.
Estas conclusiones cierran el caso de estudio completo y no únicamente la prueba del 20 de agosto
de 2026, porque cada decisión de interfaz que esa prueba puso a examen procede de un artefacto
anterior y se rastrea hasta él.

\subsection{El arco del caso de estudio, de la dispersión del dato a la medición con usuarios}

El punto de partida fue un problema institucional, no una idea de producto: las fuentes de
crédito, liquidez, derivados y otros dominios se consultan en sistemas separados, con
definiciones, permisos y formatos que no siempre son comparables. La Actividad 1 convirtió esa
dispersión en ocho personas con sus mapas de empatía y eligió a Laura Méndez como persona
primaria, siguiendo el criterio de que un producto se diseña para un blanco específico y no para
un promedio de usuarios (Hartson y Pyla, 2012). La Actividad 2 tradujo esas necesidades en seis
escenarios y cuatro journey maps, y fijó el principio que sostiene todo lo demás: la procedencia
del dato se conserva durante el recorrido entero y no se añade al final. La Actividad 3 comparó
seis alternativas del mercado y de la práctica actual, ubicó la ventaja del producto en la
intersección de catálogo gobernado, linaje visible y consulta en lenguaje natural, y corrigió con
un card sorting de 35 tarjetas la ubicación de tableros e indicadores antes de dibujar una sola
pantalla. La Actividad 4 materializó esas decisiones en siete interfaces del producto construido y
en un sistema visual con reglas verificables. La Actividad 5 las midió.

\begin{uxtablalarga}{C{1.1cm}L{3.2cm}Z{1.0}Z{1.0}}{Recorrido del caso de estudio: qué produjo cada actividad y qué hizo la prueba con esa decisión}{\thd{Etapa} & \thd{Artefacto principal} & \thd{Decisión que fijó} & \thd{Qué hizo la prueba con ella}}
A1 & Ocho personas y ocho mapas de empatía; Laura Méndez como persona primaria & El blanco de diseño es quien valida una cifra y responde con su fuente & Los cinco participantes validan o analizan información en su trabajo: la medición cayó sobre el perfil que la Actividad 1 declaró \\
A2 & Seis escenarios y cuatro journey maps; quince actividades en siete etapas & La procedencia se conserva durante todo el recorrido y no se agrega al entregar & Confirmó dos de sus predicciones de fricción y amplió el alcance del recorrido analítico \\
A3 & Seis alternativas comparadas, card sorting de 35 tarjetas con ocho evaluadores prototipo y arquitectura revisada & Tableros e indicadores pertenecen a exploración; la ventaja está en la intersección de catálogo, linaje y consulta & Ningún participante pidió ayuda para llegar al tablero: la ubicación corregida se sostuvo ante personas reales \\
A4 & Siete interfaces del producto construido, guía de estilos y prevalidación con evaluadores sintéticos & El perfil decide los valores por omisión; toda cifra del asistente cita la consulta que la produjo & El material de la prueba fue ese producto y no una maqueta; la trazabilidad del asistente resultó la fortaleza mejor reconocida \\
A5 & Prueba con cinco participantes reales el 20 de agosto de 2026 & La corrección se ordena por severidad primero y por frecuencia después & Cerró el ciclo y dejó la línea base contra la que se compara el retest \\
\end{uxtablalarga}

La cadena de evaluación también quedó ordenada. La Actividad 4 gastó primero una ronda barata con
ocho evaluadores sintéticos condicionados por las personas de la Actividad 1 (Bougie et al.,
2026), y solo después el proyecto gastó la única ronda con participantes humanos de que disponía.
Esa secuencia, evaluación simulada antes que humana, es la que la literatura reciente describe al
automatizar la evaluación de experiencia sobre interfaces web reales (Tan et al., 2026). Su
resultado aquí es comprobable: de las dos preguntas que la
prevalidación dejó abiertas, la prueba cerró una y dejó la otra sin medir.

\subsection{El cierre del círculo: la fricción apareció donde el recorrido la había señalado}

El argumento más fuerte de este documento no es la puntuación SUS de 89.0. Es que las dos
fricciones que la prueba encontró estaban escritas, con otro vocabulario, en el journey map que el
equipo construyó en taller varias semanas antes de tener usuarios enfrente.

\textbf{Dónde buscar.} El listado de actividades de la Actividad 2 descompuso el objetivo de Laura
Méndez en quince pasos, y el segundo de todos ---dentro de la etapa de disparo, antes de tocar el
portal--- era \emph{decidir dónde buscar}: elegir entre el portal, un archivo anterior o preguntar
a un colega. La tabla de brechas de conocimiento de esa misma actividad registró para esa etapa la
pregunta que Laura no sabe responder: si el portal contiene el dato que le piden. Esa etapa
produjo la fricción más frecuente del estudio. H2 apareció en tres de los cinco participantes: C y
V probaron primero el buscador de cabecera y A1 dudó qué parte utilizar primero. La sub-escala de
aprendizabilidad corrobora el hallazgo con un segundo instrumento: A1, el participante que
titubeó, obtuvo 75.0 en aprendizabilidad frente a 84.4 en usabilidad, la única separación
apreciable entre las dos sub-escalas de la muestra. La predicción se confirmó en su naturaleza y
se desplazó de lugar: la indecisión no se resolvió al entrar al portal, sino que sobrevivió dentro
de él como duda entre dos controles de búsqueda. R2 atiende exactamente ese desplazamiento.

\textbf{Confiar en la proyección.} El escenario del perfil directivo registró como punto de dolor
que en una reunión un director presente un número y otro presente uno distinto porque usaron
cortes de fecha diferentes, y respondió con el sello de conciliación y la fecha de corte visibles,
\emph{de manera que la persona sabe con certeza qué versión está presentando}. La prueba confirmó
la categoría del riesgo y precisó su forma: la confusión no ocurrió entre dos cortes, sino entre
la proyección y el último dato observado dentro de una misma serie. V interpretó primero la
proyección como último dato; T4 fue el único no-éxito del estudio, con 2 errores, 94 segundos y
SEQ 4, contra un SEQ promedio de 5.8 de 7. El hallazgo aparece en un solo participante y encabeza
aun así la lista de correcciones, porque la severidad por impacto y la frecuencia son dos
criterios distintos y el primero manda cuando la consecuencia es presentar una cifra equivocada en
una junta directiva (Albert y Tullis, 2013, cap.~5).

\begin{uxtablalarga}{L{2.6cm}Z{1.0}Z{1.1}L{2.3cm}}{Predicciones de las Actividades 1 a 4 frente a lo observado el 20 de agosto de 2026}{\thd{Origen} & \thd{Predicción} & \thd{Qué encontró la prueba} & \thd{Veredicto}}
A2, journey, etapa de disparo & Decidir dónde buscar es una decisión previa al portal & H2 en 3 de 5: C y V probaron el buscador de cabecera y A1 dudó qué parte usar primero & Confirmada y desplazada al interior del portal \\
A2, escenario directivo & Un indicador sin versión visible expone a presentar la cifra equivocada & H1: V leyó la proyección como último dato; T4 parcial, 2 errores, SEQ 4 & Confirmada, con una forma que el escenario no nombró \\
A2, curva emocional & La validación es el momento de la verdad y el punto más alto del recorrido con el portal & T2 con éxito en C y V y SEQ 6/6; H6 como fortaleza en los tres participantes expuestos & Confirmada \\
A2 y A4, tarjeta de llamada a herramienta & Hacer visible la consulta y la fuente sostiene la confianza en la cifra & H5: C y V identificaron consulta y fuente sin ayuda, con SEQ 6/6 en T5 y cero solicitudes de ayuda & Confirmada \\
A3, prueba de árbol y card sorting & Tableros e indicadores pertenecen a exploración: en la ubicación original solo 1 de 8 evaluadores acertó la ruta y 6 de 8 los buscaron en exploración & Nadie pidió ayuda para llegar al tablero y T4 se resolvió dentro de su umbral de 120 s, en 76 y 94 segundos; la fricción restante fue de interpretación & Confirmada; el riesgo pendiente cambió de naturaleza \\
A4, prevalidación sintética & La bitácora de accesos conserva un momento de decisión entre administración y gobierno, la única tarea con duda unánime de los ocho evaluadores & Ninguna de las tareas del 20 de agosto recorrió esa ruta & Sin medir \\
A4, valores por omisión por perfil & El perfil decide la configuración y el usuario no la ajusta & H4: V buscó Administración antes de cambiar de perfil y volvió a revisar la cabecera; la regla se sostiene y la transición necesita anuncio & Matizada \\
A2, alcance de los seis escenarios & El recorrido analítico termina en la extracción del archivo & H7: A2 pidió agregaciones o transformaciones sencillas antes de extraer & Ampliada: el recorrido tiene un paso que los escenarios no modelaron \\
\end{uxtablalarga}

Tres consecuencias se desprenden de la tabla. La primera es que el journey map funcionó como
instrumento de predicción y no solo como pieza de documentación: señaló dos etapas y las dos
resultaron ser las que fallaron. La segunda es que la prevalidación con evaluadores sintéticos
acertó en la corrección estructural ---la nueva ubicación de los tableros se sostuvo con personas
reales--- pero no anticipó la fricción interpretativa, que es del tipo que requiere a una persona
mirando una gráfica. La tercera es la única ampliación del modelo: los escenarios de la Actividad
2 cerraban el recorrido analítico en la extracción, y un participante lo continuó un paso más allá
pidiendo transformaciones ligeras antes de extraer. Ese es el hallazgo que el trabajo de las
semanas anteriores no contenía.

El diseño de dos bloques ---tareas estandarizadas para C y V, recorridos vinculados con el trabajo
para A1, A2 y A3--- fue el que permitió que ambas cosas ocurrieran en la misma sesión: el primer
bloque produjo métricas comparables y el segundo, relevancia y descubribilidad en contexto real.
El mismo reparto entre escenario controlado y tarea propia del participante aparece en las
plataformas actuales de evaluación de la interacción con agentes conversacionales, como HARP
(Zhu et al., 2026), y aquí resultó decisivo: H7 y H8 solo pudieron aparecer en el bloque por
contexto, y H1 solo pudo cuantificarse en el estandarizado.

\subsection{La siguiente iteración}

Las siete recomendaciones quedan ordenadas por dos criterios y no por uno. La severidad por
impacto encabeza la lista y la frecuencia ordena el segundo tramo, que es la combinación que la
literatura de métricas recomienda para priorizar problemas de usabilidad (Albert y Tullis, 2013,
cap.~5). El resultado es este orden de trabajo.

\begin{uxtablalarga}{C{0.9cm}L{3.4cm}L{2.2cm}Z{1.0}}{Orden de la siguiente iteración: las siete recomendaciones con su prioridad y su criterio de cierre}{\thd{ID} & \thd{Recomendación} & \thd{Prioridad} & \thd{Qué la da por cerrada}}
R1 & Separar observado y proyectado con rótulo persistente, fecha de corte y lenguaje inequívoco & 1, crítica por impacto & Repetir T4 y obtener interpretación correcta sin ayuda, con el SEQ del tablero por encima del 4.5 actual \\
R2 & Unificar la intención de búsqueda: distinguir la búsqueda global de la contextual y conservar el término & 2, alta por frecuencia, 3 de 5 & Repetir T2 observando dónde arranca la búsqueda, sin cambio de control a mitad de tarea \\
R3 & Llevar los límites al selector de exportación: tamaño, compatibilidad y uso recomendado junto a cada formato & 3, media & Repetir T3 con las restricciones consultadas antes de elegir formato y no después \\
R4 & Confirmar el contexto al cambiar de perfil, con redirección al inicio del perfil o aviso con acceso directo & 4, media & Repetir T6 sin revisión adicional de la cabecera tras el cambio \\
R5 & Conservar visibles los metadatos, la fuente, el responsable, la consulta y el estado & Restricción permanente & Verificar que R1 a R4 no degraden H5 ni H6 en el retest \\
R6 & Evaluar transformaciones ligeras antes de comprometer desarrollo & Investigación previa & Observar a más analistas y decidir con esa evidencia, no con la señal de un participante \\
R7 & Retestar contra esta línea base & Cierre del ciclo & Repetir T2 y T4 y cronometrar cada actividad del bloque por contexto \\
\end{uxtablalarga}

El retest no es un trámite de cierre: es lo que convierte una medición aislada en un ciclo de
medir, corregir y volver a medir. El precedente de dominio más cercano ---la evaluación de
usabilidad de una herramienta de analítica de autoservicio (Joarder et al., 2026)---
sigue esa misma secuencia y obtiene su valor de la comparación entre rondas, no de la primera. Por
eso la línea base queda escrita con precisión: 11 éxitos de 12 tareas, 91.7\,\%; tiempo medio de
59.5 segundos; 7 errores; cero solicitudes de ayuda; SEQ promedio de 5.8 de 7 con el tablero en
4.5; facilidad media de 6.0 de 7 en los recorridos por contexto; y SUS de 89.0 con intervalo de
confianza del 95\,\% entre 79.8 y 98.2. El límite inferior de ese intervalo supera la meta de 75 y
queda por encima del promedio de referencia de 68 que reportan los estudios acumulados de la
escala (Sauro y Lewis, 2016), de modo que la comparación posterior tiene contra qué medirse.

El plan de la siguiente ronda incorpora además dos ausencias declaradas de esta. La primera son
los perfiles principiantes: los cinco participantes trabajan con información y la muestra, elegida
así a propósito, no dice nada sobre quien se acerca al portal sin ese oficio. La segunda es la
bitácora de accesos, la ruta que hizo titubear a los ocho evaluadores sintéticos y que ninguna
tarea del 20 de agosto recorrió. Ambas quedan como preguntas abiertas con su instrumento ya
elegido, no como omisiones.

\subsection{Lo que el caso de estudio deja escrito}

La conclusión metodológica del proyecto es que Karisma Data no se mide por la cantidad de fuentes
que reúne. Su éxito depende de que una persona encuentre el dato correcto, comprenda su contexto,
complete la tarea sin perder la procedencia y reconozca los límites de lo que el sistema le
muestra. Las tres fortalezas que la prueba confirmó ---la trazabilidad del asistente, los
metadatos que permiten reconocer una fuente y una arquitectura donde nadie se perdió---
pertenecen a esa definición, y las dos fricciones que encontró también: una es no saber por dónde
empezar y la otra es creer que un número es lo que no es.

El ciclo de análisis, diseño, prototipo y evaluación que describe la literatura de proceso
(Hartson y Pyla, 2019) quedó recorrido una vez completo, con evidencia propia en cada etapa y sin
saltarse la última, que es la que suele quedar fuera de un trabajo académico por falta de acceso a
usuarios. El material evaluado tampoco fue una maqueta: fue el producto construido, y
\siHayDemoWeb{la misma versión que recorrieron los cinco participantes se consulta en
\urlDemoWeb}{la dirección pública del prototipo se entrega junto con este documento}. Eso permite
que cualquier lector repita las seis tareas del bloque estandarizado y contraste sus resultados
con la línea base publicada aquí, que es la forma más simple de comprobar que las cifras de estas
páginas no son una afirmación del equipo sino una medición reproducible.

En la preparación editorial de este volumen el equipo se apoyó en Claude Code (Anthropic, 2026)
como herramienta de asistencia para la composición en LaTeX, la verificación cruzada de cifras
entre capítulos y la revisión de maquetación. Las decisiones de diseño, la conducción de la prueba
y la interpretación de los resultados son del equipo.

% ============================================================================
\clearpage
\section{Referencias bibliográficas}\label{sec:a5-referencias}

Esta es la lista consolidada del documento: reúne las fuentes que sostienen afirmaciones concretas
en las cinco partes, en formato APA y con sangría francesa. No incluye lecturas que no se citen en
el cuerpo.

\begin{uxreferencias}
Albert, B., y Tullis, T. (2013). \textit{Measuring the user experience: Collecting, analyzing, and
presenting usability metrics} (2.\textsuperscript{a} ed.). Morgan Kaufmann.

Annacchino, M. A. (2003). \textit{New product development: From initial idea to product
management}. Elsevier.

Anthropic. (2026). \textit{Claude Code} (modelo Claude Fable 5) [Software de asistencia basado en
modelos de lenguaje]. \href{https://claude.com/claude-code}{claude.com/claude-code}.

Atlan. (2026). \textit{Atlan named a Leader in the 2026 Gartner Magic Quadrant for D\&A Governance
Platforms}. Consultado el 9 de agosto de 2026 en
\href{https://atlan.com/}{sitio oficial de Atlan}.

Baxter, K., Courage, C., y Caine, K. (2015). \textit{Understanding your users: A practical guide
to user research methods} (2.\textsuperscript{a} ed.). Morgan Kaufmann.

Bloomberg L.P. (2026). \textit{Bloomberg Professional Services}. Consultado el 6 de agosto de 2026
en \href{https://www.bloomberg.com/professional/}{sitio oficial de Bloomberg Professional}.

Bougie, N., Ye, X., Marconi, G. M., y Watanabe, N. (2026). \textit{PerceptUI: LLM agents as
human-aligned synthetic users for UI/UX evaluation}. arXiv.
\href{https://arxiv.org/abs/2606.05697}{arxiv.org/abs/2606.05697}.

Brooke, J. (1996). SUS: A quick and dirty usability scale. En P. W. Jordan, B. Thomas,
B. A. Weerdmeester y A. L. McClelland (Eds.), \textit{Usability evaluation in industry}
(pp. 189--194). Taylor \& Francis.

Collibra. (2026). \textit{Collibra Data Intelligence Platform}. Consultado el 9 de agosto de 2026
en \href{https://www.collibra.com/}{sitio oficial de Collibra}.

de Voil, N. (2020). \textit{User experience foundations}. BCS, The Chartered Institute for IT.

Hartson, R., y Pyla, P. S. (2012). \textit{The UX book: Process and guidelines for ensuring a
quality user experience}. Morgan Kaufmann.

Hartson, R., y Pyla, P. S. (2019). \textit{The UX book: Agile UX design for a quality user
experience} (2.\textsuperscript{a} ed.). Morgan Kaufmann.

International Organization for Standardization. (2018). \textit{Ergonomics of human-system
interaction: Part 11: Usability: Definitions and concepts} (ISO 9241-11:2018).

Jang, J., y Li, W.-S. (2026). \textit{TwinBI: An agentic digital twin for efficient augmented
interactions with business intelligence dashboards}. arXiv.
\href{https://arxiv.org/abs/2606.13731}{arxiv.org/abs/2606.13731}.

Joarder, S., Chatti, M. A., y Born, L. (2026). \textit{Usability evaluation and improvement of a
tool for self-service learning analytics}. arXiv.
\href{https://arxiv.org/abs/2603.24321}{arxiv.org/abs/2603.24321}.

LSEG. (2026). \textit{LSEG Workspace}. Consultado el 9 de agosto de 2026 en
\href{https://www.lseg.com/}{sitio oficial de LSEG}.

Maioli, L. (2018). \textit{Fixing bad UX designs: Master proven approaches, tools, and techniques
to make your user experience great again}. Packt Publishing.

McInerney, P. (2003). Plantilla de escenarios [citado en Baxter, K., Courage, C., y Caine, K.
(2015), \textit{Understanding your users: A practical guide to user research methods},
2.\textsuperscript{a} ed., Morgan Kaufmann].

Microsoft. (2024). \textit{Important update to Microsoft Power BI pricing}. Consultado el 6 de
agosto de 2026 en
\href{https://powerbi.microsoft.com/blog/important-update-to-microsoft-power-bi-pricing/}{blog
oficial de Power BI}.

Microsoft. (2026a). \textit{Deprecating Power BI Q\&A}. Consultado el 9 de agosto de 2026 en
\href{https://powerbi.microsoft.com/en-us/blog/deprecating-power-bi-qa/}{blog oficial de Power BI}.

Microsoft. (2026b). \textit{Microsoft Purview pricing}. Consultado el 9 de agosto de 2026 en
\href{https://azure.microsoft.com/en-us/pricing/details/purview/}{página oficial de precios de
Purview}.

Morville, P., y Rosenfeld, L. (2006). \textit{Information architecture for the World Wide Web:
Designing large-scale web sites} (3.\textsuperscript{a} ed.). O'Reilly Media.

Naik, S., Passi, S., Vorvoreanu, M., Saponas, S., y Hall, A. (2026). ``So there's a Catch-22
here'': How early adopters who build multi-agent LLM systems conceptualize transparency. arXiv.
\href{https://arxiv.org/abs/2606.08323}{arxiv.org/abs/2606.08323}.

Osterwalder, A., y Pigneur, Y. (2010). \textit{Business model generation: A handbook for
visionaries, game changers, and challengers}. John Wiley \& Sons.

Peng, Y.-H., Das, S., Bigham, J. P., y Wu, J. (2026). \textit{Efficient personalization of
generative user interfaces}. arXiv.
\href{https://arxiv.org/abs/2604.09876}{arxiv.org/abs/2604.09876}.

Portigal, S. (2013). \textit{Interviewing users}. Rosenfeld Media.

Pyramid Analytics. (2026). \textit{Decision Intelligence Platform y opciones de precios}.
Consultado el 6 de agosto de 2026 en
\href{https://www.pyramidanalytics.com/decision-intelligence-platform/get-pricing/}{página oficial
de precios}.

Ramírez Mejía, A. I. (2023). \textit{Grocery shopping app: A guided UX case study} (partes 1 a 5)
[Diapositivas]. TC4032, Experiencia del usuario y diseño de interfaces, Maestría en Inteligencia
Artificial Aplicada, Instituto Tecnológico y de Estudios Superiores de Monterrey.

Sauro, J., y Lewis, J. R. (2016). \textit{Quantifying the user experience: Practical statistics
for user research} (2.\textsuperscript{a} ed.). Morgan Kaufmann.

ServiceNow. (2026). \textit{Pyramid Analytics: Expanding data-driven AI}. Consultado el 9 de
agosto de 2026 en
\href{https://www.servicenow.com/workflow/news/pyramid-analytics-expanding-data-driven-ai.html}{sala
de prensa de ServiceNow}.

Spool, J. (2005). \textit{The anatomy of a design decision}. UIE.

Tan, W. J., Lim, Z. R. L., Durgad, S., Obegi, K., y Li, A. Y. (2026). \textit{Avenir-UX: Automated
UX evaluation via simulated human web interaction with GUI grounding}. arXiv.
\href{https://arxiv.org/abs/2604.09581}{arxiv.org/abs/2604.09581}.

ThoughtSpot. (2026). \textit{ThoughtSpot plans and pricing}. Consultado el 9 de agosto de 2026 en
\href{https://www.thoughtspot.com/pricing}{página oficial de precios de ThoughtSpot}.

Tullis, T., y Wood, L. (2004). How many users are enough for a card-sorting study?
\textit{Proceedings of the Usability Professionals Association Conference}.

Velasco, J., Morales, L., y Penado, C. (2022). Arquitectura de la información. En M. Del Río y
F. Linares (Eds.), \textit{UX Latam: Historias sobre definición y diseño de servicios digitales}
(pp. 47--63). Universidad del Pacífico.

W3C. (2023). \textit{Web Content Accessibility Guidelines (WCAG) 2.2}. World Wide Web Consortium.
Consultado el 10 de agosto de 2026 en \href{https://www.w3.org/TR/WCAG22/}{w3.org/TR/WCAG22/}.

Walter, S. (2022). \textit{User journey mapping: Visualize user research, brainstorm
opportunities, and solve problems}. SitePoint.

Zhu, Z., Friedman, N., Weatherwax, K., y Eiben, E. (2026). \textit{HARP: The Human-AI Research
Platform}. arXiv. \href{https://arxiv.org/abs/2607.20773}{arxiv.org/abs/2607.20773}.
\end{uxreferencias}

% ============================================================================
\clearpage
\section{Anexo A. Instrumentos de aplicación}

Este anexo reúne los instrumentos tal como se aplicaron el 20 de agosto de 2026. El manual de uso
del prototipo se publica aparte, como Anexo B.

\subsection{Valoración general de facilidad}

Al terminar las actividades se preguntó: ``En una escala de 1, muy difícil, a 7, muy fácil,
¿qué tan fácil fue realizar las actividades asignadas?''. Se añadieron dos preguntas abiertas:
``¿qué parte te hizo detenerte o dudar?'' y ``¿qué cambiarías primero?''.

\subsection{Cuestionario SUS}

Cada afirmación se respondió de 1, totalmente en desacuerdo, a 5, totalmente de acuerdo. Los diez
reactivos son los de la escala original (Brooke, 1996), en su redacción en español.

\begin{uxpreg}
  \item Creo que me gustaría utilizar este sistema con frecuencia.
  \item Encontré el sistema innecesariamente complejo.
  \item Pensé que el sistema era fácil de usar.
  \item Creo que necesitaría apoyo de una persona con conocimientos técnicos para utilizar este sistema.
  \item Encontré que las distintas funciones del sistema estaban bien integradas.
  \item Pensé que había demasiada inconsistencia en el sistema.
  \item Imagino que la mayoría de las personas aprendería a utilizar este sistema rápidamente.
  \item Encontré el sistema muy incómodo o engorroso de utilizar.
  \item Me sentí seguro al utilizar el sistema.
  \item Necesité aprender muchas cosas antes de poder utilizarlo.
\end{uxpreg}

\subsection{Campos de la hoja de observación}

Código de participante, área, años de experiencia, modalidad, actividad asignada, hora de inicio,
duración total, resultado, solicitud de ayuda, duda o retroceso, comentario del participante,
facilidad general, respuestas SUS y observación de la persona moderadora. El consentimiento y las
notas identificables se conservan fuera del entregable durante el periodo informado.
